\section{Public-Data Audits}
\label{sec:data}

\paragraph{Releases and distinct validation targets.}
Chicago's TNP release supplies a declared two-transaction shared-trip
benchmark with nearest-15-minute times and coarsened geography
\citep{chicagodata2026,chicagomanual2026,chicagoprivacy2026}.
NYC's 2023 High Volume For-Hire Vehicle table supplies second-level pickup and
drop-off times, shared-match indicators, Taxi Zones, miles, and trip seconds,
but no public event key. Neither release reveals operational co-membership.
Chicago therefore tests extraction and optimization on a public-envelope
$K=2$ graph; NYC tests fixed-time ordered-event optimization. These are
different validation targets: success on the graph benchmark does not verify
joint timestamp realizability or identify actual vehicle runs.

\paragraph{Frozen cohorts and certification.}
Candidate rules use public temporal information and boundary roles; spatial
screening and omitted-incidence budgets are explicit sensitivity parameters.
The Chicago calendar and eligibility rules precede each campaign's outcome
evaluation, and ineligible dates are retained without replacement. We report
the original pilot separately from the calendar-disjoint follow-up, which was
declared after inspecting the pilot. Neither is a probability sample or a
pristine holdout. Repeated queries and parameter settings within one window
are not independent cohorts. A numerical endpoint certificate requires an
optimal solver status and witness replay under the recorded tolerances;
this is distinct from exact-arithmetic verification. Missing public outcomes,
transport failures, and unresolved optimization have separate denominators.
Published artifacts contain aggregate results and provenance hashes rather
than row identifiers, partner assignments, or latent timestamps.

\section{Results and Their Boundaries}
\label{sec:results}

\subsection{Chicago: graph optimization closes before the study does}
The pilot predeclares 24 weekday/time windows from January 5--14, 2026.
All return terminal records: 23 complete and one scientifically ineligible.
The completed windows contain 34--77 core rows and 363--913 candidate rows.
Every completed extraction passes public source-count closure, which verifies
the declared temporal extraction, not true-partner completeness.
Table~\ref{tab:chicagopanel} keeps the query and window denominators separate.
All 2,744 pairs with complete public query values close numerically; the other
686 pairs lack required public values and publish no certified interval.

\begin{table}[t]
\caption{Chicago public-envelope $K=2$ audits. Endpoint rows count query--parameter
pairs in completed windows. Complete-data rate excludes missing values
but includes computationally unresolved pairs. Follow-up is an interim snapshot.}
\label{tab:chicagopanel}
\centering
\small
\begin{tabular}{@{}lrr@{}}
\toprule
Quantity & Pilot & Follow-up \\
\midrule
Declared windows & 24 & 96 \\
Completed / ineligible & 23 / 1 & 10 / 2 \\
Execution failures / unstarted & 0 / 0 & 0 / 84 \\
Endpoint pairs & 3,430 & 1,500 \\
Numerically certified pairs & 2,744 & 1,192 \\
Missing-public-value pairs & 686 & 300 \\
Computationally unresolved pairs & 0 & 8 \\
Complete-data rate & 100\% & 99.33\% \\
\bottomrule
\end{tabular}
\end{table}

The follow-up predeclares 96 windows from January 15--February 27 under the
same eligibility and 60-second endpoint budgets. Its first 12 windows produce
10 completed records and two ineligible records; 84 windows remain unstarted.
The eight unresolved pairs all concern duration-gap lower endpoints in the
January 15 evening window. Their feasible incumbents are not certified
minima, so the pair-level intervals remain withheld. Thus a completed window
does not mean every query is certified, and 120 declared windows across the
two designs do not mean 120 completed or eligible cohorts.

\paragraph{Stable support sensitivity can coexist with weak identification.}
One frozen January 13 evening audit illustrates the distinction. Completing
the boundary adds 92 buffers and 2,629 edges, increasing the graph by 12.15\%.
Miles-gap width grows only 0.53\%, from 20.159 to 20.266 miles; duration-gap
width grows 1.02\%, from 46.470 to 46.942 minutes. The boundary correction
therefore has little effect on these already broad intervals. Community-area
fractions retain the full $[0,1]$ range. Moreover, 48.88\% of temporal edges
lack complete public endpoint centroids and remain in every radius screen.
Small radius sensitivity consequently does not establish geographic
precision. These are descriptive results for one declared candidate graph,
not evidence of operational co-rider identity or citywide prevalence.

\subsection{NYC: decision ambiguity need not await optimal endpoints}
The frozen exact-second decision panel declares 24 windows: 20 eight-core
windows across seasons, weekday/weekend regimes, and times of day, plus four
16-core stress windows. Twenty-one are eligible and three fail the
outcome-blind core-integrity/cap criteria; technical failures are zero.
Two selected-buffer outcomes, miles and trip minutes, at $C=2,3,4$ yield 126
cells (Table~\ref{tab:decisionpanel}). Of these, 101 endpoint pairs close;
125 cells have replayed feasible values on both sides of the candidate-median
threshold, certifying ambiguity even where extrema remain unresolved.
One 16-core miles cell has an unresolved decision status.

\begin{table}[t]
\caption{NYC decision panel. Ambiguity requires attained values on both sides;
point disagreement compares deterministic methods, not predictions with truth.}
\label{tab:decisionpanel}
\centering
\scriptsize
\setlength{\tabcolsep}{2.5pt}
\begin{tabular}{@{}crrrrrr@{}}
\toprule
Core & Eligible & Cells & Closed & Ambiguous & Unresolved & Point split \\
\midrule
8 & 17 & 102 & 91 & 102 & 0 & 24 \\
16 & 4 & 24 & 10 & 23 & 1 & 0 \\
All & 21 & 126 & 101 & 125 & 1 & 24 \\
\bottomrule
\end{tabular}
\end{table}

Nearest-time and zone-aware pairings, together with temporal-score and
zone-score ordered decompositions, disagree in 24/126 cells (19.0\%).
All 102 eight-core cells remain ambiguous at the candidate first quartile,
median, and third quartile. These thresholds probe relation sensitivity;
they are not validated operational decision thresholds. Because event truth
is absent, neither ambiguity nor point disagreement measures prediction error.

\subsection{The structural comparator returns an informative null}
A separate retrospective four-core/twelve-buffer audit compares ordered
events, pairs, and overlap cliques on identical exact times, roles, and
outcomes. Every positive selected-buffer count $q$ feasible in all three
models is compared, holding the estimand fixed. Across two outcomes, three
capacities, two common counts, and two restrictions, \emph{none of the 24
endpoint-pair comparisons changes}. The reason is observable: all 16
intervals share 14 seconds, yielding a complete overlap graph. Every
capacity-feasible ordered event is already a clique on this input.

At $C=3,4$, ordered events reach eight and twelve selected buffers, whereas
pairs reach four. Those differences concern attainable support; they do not
demonstrate a changed outcome at a common $q$. This audit therefore provides
no empirical advantage for sequential turnover over clique modeling.
Its result is retained rather than replacing the cohort after inspection.

A subsequent geometry census retains all 24 declared cells. Recovery leaves
eight geometry-complete, three scientifically ineligible, and thirteen
transport-unresolved cells (eleven failures and two deadlines). All eight
completed full views contain a core-touching non-clique event column, while
four of their time-only $4+12$ views are cliques. A local non-clique column
does not establish a complete feasible world or a fixed-$q$ outcome effect.
The incomplete census therefore cannot resolve the structural audit's null
or establish population prevalence.

\subsection{Integer support optimization has a bounded scale claim}
The separate branch-and-price lattice closes all 18 maximum-support problems
at six sizes from $4+12$ through $16+48$ and $C=2,3,4$
(Appendix Table~\ref{tab:scaledetail}). Four initially open cells close under
longer budgets; their earlier anytime bounds remain in the audit history.
This establishes certified support optimization on one deterministic cohort.
It does not close the 25 numerically unclosed outcome pairs in the decision panel, provide
a city-scale runtime guarantee, or recover true memberships. The evidence
supports conditional computation and substantial relation ambiguity; the
empirical importance of the full temporal-event structure remains open.
